\section*{第 \ref{ch:g10:chemistry-of-life} 章——生物体的化学组成}

\begin{solution}{exo:g10:chemistry-of-life:1}
氧（约 65\%）、碳（18.5\%）、氢（9.5\%）和氮
（3.2\%）：四者合计占体重的 96\%。
\end{solution}

\begin{solution}{exo:g10:chemistry-of-life:2}
\omterm{def:g10:chemistry-of-life:mineral-organic}{无机物}：水、氯化钠、磷酸钙、二氧化碳（一种碳化合物，
但在生物体外形成，且不含碳–氢骨架）。
\omterm{def:g10:chemistry-of-life:mineral-organic}{有机物}：葡萄糖、甘油三酯、\omterm{def:g10:chemistry-of-life:families}{氨基酸}。
\end{solution}

\begin{solution}{exo:g10:chemistry-of-life:3}
失水 $12.0 - 0.6 = \qty{11.4}{g}$；$11.4/12.0 = 0.95$，可知叶片
含水 95\%。
\end{solution}

\begin{solution}{exo:g10:chemistry-of-life:4}
(a) 碘液滴在捣碎的种子上：蓝黑色。(b) 双缩脲
试剂：紫色。(c) 将种子在纸上捣碎：形成一个不干的半透明斑点，
不会变干；苏丹 III 将其染成红色。
\end{solution}

\begin{solution}{exo:g10:chemistry-of-life:5}
\omterm{def:g10:chemistry-of-life:families}{糖类}：葡萄糖等单糖。\omterm{def:g10:chemistry-of-life:families}{脂质}：甘油和
脂肪酸。\omterm{def:g10:chemistry-of-life:families}{蛋白质}：\omterm{def:g10:chemistry-of-life:families}{氨基酸}（二十种）。\omterm{def:g10:chemistry-of-life:families}{核酸}：
核苷酸（四种）。
\end{solution}

\begin{solution}{exo:g10:chemistry-of-life:6}
碳：浓集 $18.5/0.02 \approx 900$ 倍。氮：
$3.2/0.002 = 1600$ 倍。生物物质并非其外界的样本：
它将若干稀有元素富集了上千倍，
这正是生命的化学特征。
\end{solution}

\begin{solution}{exo:g10:chemistry-of-life:7}
水：$0.62 \times 70 = \qty{43.4}{kg}$。干物质：$70 - 43.4 =
\qty{26.6}{kg}$。\omterm{def:g10:chemistry-of-life:families}{蛋白质}：$0.17 \times 70 = \qty{11.9}{kg}$。
\end{solution}

\begin{solution}{exo:g10:chemistry-of-life:8}
\omterm{def:g10:chemistry-of-life:families}{糖类} $60 \times 17 = \qty{1020}{kJ}$；\omterm{def:g10:chemistry-of-life:families}{脂质} $12 \times 38 =
\qty{456}{kJ}$；\omterm{def:g10:chemistry-of-life:families}{蛋白质} $8 \times 17 = \qty{136}{kJ}$；总计
\qty{1612}{kJ}，约 \qty{1600}{kJ}。\omterm{def:g10:chemistry-of-life:families}{脂质}提供 $456/1612 \approx
28\%$ 的能量，而其质量仅占 12\%。
\end{solution}

\begin{solution}{exo:g10:chemistry-of-life:9}
面包：呈蓝黑色——面粉主要成分是淀粉。奶酪：除碘液自身的颜色外无变化，
——含\omterm{def:g10:chemistry-of-life:families}{蛋白质}和\omterm{def:g10:chemistry-of-life:families}{脂质}，不含淀粉。糖：
无变化——蔗糖是单糖，而非淀粉；碘只能检测
长链淀粉。
\end{solution}

\begin{solution}{exo:g10:chemistry-of-life:10}
干物质主要是\omterm{def:g10:chemistry-of-life:mineral-organic}{有机物}：含碳–氢的分子，
能与氧结合，释放能量，生成二氧化碳和水。灰分
为\omterm{def:g10:chemistry-of-life:mineral-organic}{无机盐}，已完全与氧结合；其中
没有能够燃烧的物质。
\end{solution}

\begin{solution}{exo:g10:chemistry-of-life:11}
斐林试剂与单个葡萄糖分子所携带的游离化学基团发生反应；
在淀粉中，葡萄糖单元彼此相连，
这些基团被锁在链内，因此试剂
找不到作用对象。相反，碘会嵌入盘绕的长链淀粉中，
而单个葡萄糖分子无法形成这种链。每一种检测
检测的是结构，而非元素。
\end{solution}

\begin{solution}{exo:g10:chemistry-of-life:12}
谷粒大量吸水。含水量为 13\% 时，储备物质呈固体状态，
酶无法工作：谷粒的化学反应处于关闭状态，
以休眠状态存活。水溶解储备物质，使酶得以发挥作用，
并把产物运送到生长中的胚胎：化学反应只有在水充足时
才会重新启动。
\end{solution}

\begin{solution}{exo:g10:chemistry-of-life:13}
食盐：$9 \times 0.5 = \qty{4.5}{g}$；葡萄糖：$1 \times 0.5 =
\qty{0.5}{g}$。纯水会稀释血浆；水随后
会进入红细胞，使其膨胀破裂。
\end{solution}

\begin{solution}{exo:g10:chemistry-of-life:14}
糖原：$500 \times 17 = \qty{8500}{kJ}$。\omterm{def:g10:chemistry-of-life:families}{脂质}：$10000 \times 38 =
\qty{3.8e5}{kJ}$。若要以糖原形式储存 \qty{3.8e5}{kJ}：$380000/17
\approx \qty{22.4}{kg}$ 的糖原，再加上该质量三倍的
水，共约 \qty{90}{kg}——超过此人的体重。脂肪
每千克所携带的能量约多出九倍，这也是
需要移动的动物以\omterm{def:g10:chemistry-of-life:families}{脂质}形式储备能量的原因。
\end{solution}

\begin{solution}{exo:g10:chemistry-of-life:15}
这些原子确实普通，但（1）其比例并不普通：生物
物质将碳富集了上千倍，氮的富集程度更高，
与由氧和硅构成的地壳相比；（2）其分子也不普通：只有
生物能将碳组装成巨大的链——\omterm{def:g10:chemistry-of-life:families}{蛋白质}、淀粉、
DNA——其大小在矿物中没有对应物；（3）这些分子
以单元的排列顺序携带信息，而晶体则不具备这一特征。
生命的特殊之处在于其化学组织，而非其原子。
\end{solution}

\begin{solution}{pb:g10:chemistry-of-life:1}
\textbf{1.} 水：$0.13 \times 45 = \qty{5.9}{mg}$；干质量
\qty{39.2}{mg}。

\textbf{2.} 淀粉：$0.70 \times 39.2 \approx \qty{27.4}{mg}$；\omterm{def:g10:chemistry-of-life:families}{蛋白质}
$0.12 \times 39.2 \approx \qty{4.7}{mg}$；\omterm{def:g10:chemistry-of-life:families}{脂质} $0.02 \times 39.2
\approx \qty{0.8}{mg}$。

\textbf{3.} 碘液（蓝黑色）检测淀粉；双缩脲试剂（紫色）用于检测
\omterm{def:g10:chemistry-of-life:families}{蛋白质}；纸上出现半透明斑点或用苏丹 III 检测\omterm{def:g10:chemistry-of-life:families}{脂质}。

\textbf{4.} 含水量为 13\% 时，储备物呈固态，酶无活性；
水须先溶解储备物、使反应得以进行，
胚胎才能生长。萌发始于吸水，而
幼苗的 88\% 即说明其化学过程已经启动。

\textbf{5.} 水 $0.13 \times 8 = \qty{1.04}{t}$；干物质
\qty{6.96}{t}，约 \qty{7}{t}。

\textbf{6.} 质量 $500 + 320 + 10 - 130 = \qty{700}{g}$。面包中的
水：来自面粉 $0.13 \times 500 = \qty{65}{g}$，加上 \qty{320}{g}，
减去 \qty{130}{g}：\qty{255}{g}，即 $255/700 \approx 36\%$。

\textbf{7.} 面粉干物质 $0.87 \times 500 = \qty{435}{g}$；淀粉
$0.70 \times 435 \approx \qty{305}{g}$；\omterm{def:g10:chemistry-of-life:families}{蛋白质} $0.12 \times 435 \approx
\qty{52}{g}$（\omterm{def:g10:chemistry-of-life:families}{脂质}约 \qty{9}{g}）。

\textbf{8.} $305 \times 17 + 52 \times 17 + 9 \times 38 \approx 5185 +
884 + 342 \approx \qty{6400}{kJ}$.

\textbf{9.} 所占比例 $120/700 \approx 0.17$，约 \qty{1100}{kJ}，
即 \qty{10000}{kJ} 的 11\%。

\textbf{10.} 食盐属于\omterm{def:g10:chemistry-of-life:mineral-organic}{无机物}，已完全氧化，没有可燃烧的
碳–氢键。在体内，其离子决定血液的盐度，
并为神经与肌肉信号所必需。

\textbf{11.} 水的质量 $0.60 \times 60 = \qty{36}{kg}$。

\textbf{12.} \omterm{def:g10:chemistry-of-life:families}{蛋白质} $0.17 \times 60 = \qty{10.2}{kg}$；\omterm{def:g10:chemistry-of-life:families}{脂质} $0.16
\times 60 = \qty{9.6}{kg}$；\omterm{def:g10:chemistry-of-life:mineral-organic}{无机盐} \qty{3.6}{kg}；\omterm{def:g10:chemistry-of-life:families}{糖类}
\qty{0.6}{kg}.

\textbf{13.} 占她体内水分的 $1.2/36 \approx 3.3\%$。血容量和
体温调节都依赖于水分；失水量达 10\% 即有危险，因此
必须在数小时内通过饮水补足。

\textbf{14.} 血浆 $0.55 \times 4.5 \approx \qty{2.5}{L}$；水 $0.9
\times 2.5 \approx \qty{2.2}{kg}$，约占体内水分的 6\%。其余
位于细胞内（约占全部体内水分的三分之二）以及
细胞之间的液体中。

\textbf{15.} 婴儿每天交换的水分比例远大于
成人（包括饮水、尿液与皮肤），因此几小时内若无摄入而失水，
便会丢失其总水量中危险的比例。

\textbf{16.} \omterm{def:g10:chemistry-of-life:families}{糖类} $0.40 \times 0.6 = \qty{0.24}{kg}$；\omterm{def:g10:chemistry-of-life:families}{脂质}
$0.77 \times 9.6 \approx \qty{7.4}{kg}$；\omterm{def:g10:chemistry-of-life:families}{蛋白质} $0.53 \times 10.2
\approx \qty{5.4}{kg}$。合计约含碳 \qty{13}{kg}。

\textbf{17.} $13/60 \approx 22\%$，对比 18.5\%。两者
吻合良好，考虑到取整后的比例以及个体间\omterm{def:g10:chemistry-of-life:families}{脂质}含量的差异，
体型较瘦者得出的数值较低。

\textbf{18.} 约含氮 $0.16 \times 10.2 \approx \qty{1.6}{kg}$。它
来源于空气中的氮，经土壤细菌固定，
被植物吸收后沿食物链传递。

\textbf{19.} 氧 $0.89 \times 36 \approx \qty{32}{kg}$，即
占体重的 $32/60 \approx 53\%$。仅水分即
贡献了 65\% 中的大部分；其余为有机分子中的氧。

\textbf{20.} 一个体重 \qty{60}{kg} 的人约含
\qty{13}{kg} 碳。它原以二氧化碳形式离开空气，经
植物的光合作用固定，并经摄食与消化沿着
食物链进入人体。
\end{solution}
